\paragraph{}W ramach projektu przeprowadzono analizę wpływu modyfikacji architektury IResNet50 na skuteczność identyfikacji twarzy w warunkach okluzji. Na podstawie przeprowadzonych eksperymentów i krytycznej analizy wyników sformułowano następujące wnioski:

\begin{enumerate}
    \item \textbf{Wydajność wariantu z uwagą:} Wariant wykorzystujący moduły CBAM w blokach rezydualnych osiągnął najlepszy wynik w testach (Rank-1: 90.16\%), wykazując wysoką odporność na syntetyczną okluzję. Należy jednak zastrzec, że na sukces ten złożyła się nie tylko modyfikacja architektury, ale również zastosowanie odmiennej strategii uczenia (wyższe LR, łagodniejsza augmentacja) w porównaniu do modelu bazowego.
    \item \textbf{Wpływ gałęzi Transformer:} Wprowadzenie pomocniczej gałęzi Transformer przy zachowaniu silnej augmentacji ($p=0.7$) nie przyniosło poprawy wyników w zadanym reżimie treningowym (25 epok), co sugeruje, że bardziej złożone architektury wymagają dłuższego procesu uczenia lub większych zbiorów danych do poprawnej zbieżności.
    \item \textbf{Wyzwania metodologiczne:} Projekt uwypuklił kluczowe znaczenie rygorystycznej kontroli hiperparametrów (Learning Rate, Augmentation Probability) przy porównywaniu architektur sieci neuronowych. Zmienność tych parametrów pomiędzy eksperymentami została zidentyfikowana jako główny czynnik utrudniający jednoznaczną ocenę wpływu samych modułów uwagi.
    \item \textbf{Zalecenia do dalszych prac:} W celu weryfikacji postawionych hipotez, przyszłe prace powinny skupić się na wyrównaniu protokołu treningowego dla wszystkich modeli (strategia ceteris paribus), zastosowaniu walidacji krzyżowej oraz testach na realistycznych zbiorach okluzji (np. MAFA), aby wyeliminować wpływ czynników losowych i syntetycznego charakteru zakłóceń.
\end{enumerate}
